\section{Consequences and reusable patterns}
\label{sec:applications}

\subsection{Polynomial matrix functions}

Let $p(X)=\sum_{k=0}^d a_kX^k$ have rational coefficients. If a diagonalization certificate reconstructs $C=Q\diag(d)Q^{\mathsf T}$ with residual $\eta$, and $\bound(A,M),\bound(C,M)$ for $M\ge1$, then
\begin{equation}
 \norm{p(A)-Q\diag(p(d_i))Q^{\mathsf T}}
 \le\eta\sum_{k=1}^d k|a_k|M^{k-1}.
 \label{eq:polycalc}
\end{equation}
Indeed $p(C)=Qp(\diag(d))Q^{\mathsf T}$ by exact finite multiplication. The noncommuting telescoping identity
\[
 A^k-C^k=\sum_{j=0}^{k-1} A^j(A-C)C^{k-1-j}
\]
follows by cancellation of successive terms. Repeated left or right multiplication by $A,C$ costs at most a factor $M$ per multiplication. Thus each summand has Frobenius bound $M^{k-1}\eta$, including the endpoint cases in which one of the powers is the identity. Summing proves \eqref{eq:polycalc}. There is no erroneous Frobenius bound $\norm I=1$ in this argument: identity multiplication simply leaves the matrix unchanged.

The coefficient constant in \eqref{eq:polycalc} is a primitive recursive precision transformer. It gives stable matrix-polynomial observations despite changing diagonalization witnesses. Trace changes by at most $n$ times the Frobenius error, since each diagonal entry has at most that error. Therefore finite sums of polynomial eigenvalue values also recover the corresponding trace identities with explicit budgets.

An analytic matrix function can be handled if supplied with rational polynomial approximants, a primitive recursive uniform approximation modulus on a known interval, and the relevant finite error proof. The framework propagates those inputs. It does not manufacture the scalar approximation theorem or its modulus.

\subsection{Singular values and absolute values}

For a rational rectangular matrix $V$, its Gram matrix $A=V^{\mathsf T}V$ is positive because
\[
 x^{\mathsf T}Ax=\sum_j(Vx)_j^2\ge0.
\]
Finite elimination therefore supplies the positive certificate used above. The ordered eigenvalue names of $A$ and bounded scalar square-root searches give $n$ right-Gram singular-value names for an $m\times n$ input, with multiplicities and zeros, without any choice of singular vectors. This convention retains the additional forced zeros when $n>m$, whereas a common convention lists only $\min(m,n)$ values. The elementary scalar inequality
\[
 |\sqrt u-\sqrt v|^2\le|u-v|\qquad(u,v\ge0)
\]
is interpreted through the same rational square-root certificates; its proof is $(s-t)^2\le|s^2-t^2|$ plus the finite-slack argument.

For two finite matrices $V,W$,
\[
 V^{\mathsf T}V-W^{\mathsf T}W
 =V^{\mathsf T}(V-W)+(V-W)^{\mathsf T}W.
\]
Hence a common magnitude bound and Frobenius input error give an explicit Gram-matrix error, then Theorem~\ref{thm:weyl} gives eigenvalue errors, and the square-root comparison gives singular-value errors. This is a concrete composition of three certificate interfaces.

For symmetric represented $A$, define its matrix absolute value by
$|A|=(A^2)^{1/2}$. The preceding finite constructions supply its name without making a sign decision on any eigenvalue. The expressions $(|A|+A)/2$ and $(|A|-A)/2$ are likewise computable matrix names. By comparison with the scalar functions in arbitrarily fine diagonalization certificates, both are positive and their difference is $A$; the finite polynomial and root moduli justify this interpretation. For an explicit positivity budget, let $C=QDQ^{\mathsf T}$ approximate $A$ within $\eta$, with both Frobenius bounds at most $M$, and put $T=Q\diag(|d_i|)Q^{\mathsf T}$. Then $T^2=C^2$ and $\norm{T^2-A^2}\le2M\eta$. If $S$ is a positive spectral root certificate for $A^2$ of residual $\rho$, Theorem~\ref{thm:sqrtstable} gives
\[
 \F(S-T)\le n(\rho+2M\eta).
\]
Whenever $\sigma^2\ge n(\rho+2M\eta)$, each $(S\pm A)/2$ lies within $(\sigma+\eta)/2$ of the exactly positive matrix $Q\diag((|d_i|\pm d_i)/2)Q^{\mathsf T}$. This supplies the claimed positivity contract at every finite tolerance. A support projector or a polar factor at zero would require a separate gap contract and is not inferred from these names.

\subsection{What has been resolved}

The existing frameworks already distinguish approximate witnesses from coherent points. This report supplies a concrete method for resolving some cases where insisting on coherent witnesses would be too strong:
\begin{itemize}
\item \emph{Colliding roots in a structured subclass.} Real symmetric characteristic polynomials have ordered, multiplicity-preserving root names with a cross-input comparison theorem. This goes beyond their root supports.
\item \emph{Unstable bases.} The intrinsic square root and gap projectors have coherent names even if every finite diagonalizer selects a different basis. A global eigenbasis is not an intermediate premise.
\item \emph{Hidden finite-dimensional existence.} Rational elimination gives explicit witnesses for subspace intersection and positive certificates. Their uniform dimension bounds are part of the proof.
\item \emph{Loss of positivity during approximation.} The $A_r+e_rI$ correction and finite positivity alternative give an explicit repair with an error budget, rather than presuming that rational samples preserve a closed condition.
\item \emph{Unexplained convergence.} The spectral computation has an additive rational energy invariant and explicit fuel. The resulting name coherence follows from a separate certificate comparison theorem.
\end{itemize}

The general lesson is operational. First ask for a finite output at one tolerance. Next determine the observation that should be independent of this finite choice. Prove a comparison estimate for arbitrary certificates and distinct nearby inputs. Only then define the infinite mathematical object by its regular observations. If a discrete selection is required, isolate the numerical separation needed to make that selection stable.

\subsection{A small exact example}

For
\[
 A=\begin{pmatrix}1&1\\1&1\end{pmatrix}
\]
the first pivot has $h=0,b=1$. Taking $(c,s)=(3/5,4/5)$ gives
\[
 Q=\frac15\begin{pmatrix}3&-4\\4&3\end{pmatrix},\qquad
 Q^{\mathsf T}AQ=
 \begin{pmatrix}49/25&-7/25\\-7/25&1/25\end{pmatrix}.
\]
The original off-diagonal energy is $2$; the new energy is $98/625$. After sorting, the diagonal observations are $(1/25,49/25)$. Both orthogonality identities are exact. Further finite stages improve the residual to any prescribed rational tolerance, while the ordered observations name $(0,2)$.

The positive square root is represented without requiring the irrational eigenvector normalization to appear as an exact rational witness. Its finite certificates satisfy the explicitly checked equation $S_r^2\approx A$. At a threshold strictly between $0$ and $2$, the lower spectral projector is also a coherent matrix observation. At threshold zero there is no positive separation from the spectrum, so the same gap theorem makes no claim.
